\documentclass[../main.tex]{subfiles}
\graphicspath{{\subfix{../image/}}}

\begin{document}

\chapter{分布与叶状结构}\label{chapter:ch19-distributions-foliations}

设 $V$ 是光滑流形 $M$ 上一个处处非零的向量场。第 9 章的结果表明，$V$ 的每条积分曲线都是一个光滑浸入，并且在局部这些积分曲线的像像 Euclid 空间中的平行线那样优美地拼合在一起。流的基本定理告诉我们，这些曲线由其速度向量的知识所确定。

在本章中，我们把这一思想推广到高维子流形。一般框架如下：设 $M$ 是一个光滑流形，给定 $TM$ 的一个 $k$ 维子丛，称为 $M$ 上的一个\emph{分布}。是否存在一个 $k$ 维子流形（称为该分布的一个\emph{积分流形}），使其在每一点处的切空间都等于给定子空间？这一情形的答案比向量场情形更为复杂：分布必须满足一个非平凡的必要条件，称为\emph{对合性}。

在本章第一节中，我们定义对合性，并给出对合分布与非对合分布的例子。接着说明对合性条件如何用微分形式重新表述。

本章的主要定理，即 Frobenius 定理，告诉我们对合性也是过每一点存在积分流形的充分条件。我们以两种形式证明 Frobenius 定理。首先证明局部形式，它说每一点的一个邻域都充满了积分流形，它们像 $\mathbb{R}^n$ 中平行的仿射子空间那样优美拼合。然后证明整体形式，它说与每个对合分布相伴的是整个流形的一个划分，将其分为优美拼合的浸入积分流形，称为一个叶状结构。

在下一节中，我们应用叶状结构理论来证明 Lie 群理论中的一个重要结果。我们已经知道，对 Lie 群 $G$ 的每个 Lie 子群，都有 $\operatorname{Lie}(G)$ 的一个 Lie 子代数与之对应。利用叶状结构理论，我们证明这个对应反过来也成立：$\operatorname{Lie}(G)$ 的每个 Lie 子代数都对应于 $G$ 的某个 Lie 子群。

在本章末尾，我们给出 Frobenius 定理在偏微分方程中的若干应用。

\subfile{Chapter-19/section-01.tex}
\subfile{Chapter-19/section-02.tex}
\subfile{Chapter-19/section-03.tex}
\subfile{Chapter-19/section-04.tex}
\subfile{Chapter-19/section-05.tex}
\subfile{Chapter-19/section-06.tex}

\end{document}
